\section{Zadanie 7}
\autor{Aleksander Wieczorek}
Otrzymujemy w prezencie X kul, gdzie X ma rozkład Poissona z parametrem 2022. Dla każdej kuli rzucamy sześciościenną kostką. Jeśli wyrzuciliśmy 1, to malujemy
kulę na biało. Jeśli wyrzuciliśmy 2, to malujemy kulę na czarno. Jeśli zaś wyrzuciliśmy
coś innego, to wyrzucamy kulę do śmieci.\\
(i) Z jakim prawdopodobieństwem będziemy mieć więcej czarnych kul niż białych?\\
(ii) Jaki jest rozkład liczby czarnych kul przy założeniu, że nie uzyskaliśmy żadnej białej kuli?\\
(iii) Jaka jest oczekiwana liczba białych kul przy założeniu, że czarnych kul jest dwa razy więcej od białych?

\subsection*{Rozwiązanie}
Problem jest bardzo podobny do tego z zadania 5. Będę używał wniosków tego zadania jako blackboxa i oznaczał to gwiazdką (*)\\
Rozbijmy problem na prostszy, aby sprawnie policzyć wszystkie podpunkty. Rozbijmy najpierw zmienną X.\\
\[X' -\text{kule które zostawiamy}\]
\[W - \text{kule które wyrzucamy}\]
Zostawiamy kulę, jeśli wyrzucimy 1 lub 2 (z \(p=\frac{1}{3}\)).\\
Z (*) X' ma rozkład Poissona z parametrem \(\frac{2022}{3}=674\).
Wynik W nie będzie nas dalej interesował.\\
Analogicznie rozbijamy X':
\[C -\text{kule które pomalowano na czarno}\]
\[B - \text{kule które pomalowano na biało}\]
Zauważmy, że zostawiona kula z prawdopodobieństwem \(\frac{1}{2}\) będzie biała, albo będzie czarna z takim samym prawdopodobieństwem.\\
Jeszcze raz z (*) B oraz C są zmiennymi Poissona z parametrem \(\frac{674}{2}=337\), co więcej te zmienne są niezależne (również z (*)).
Mając to wszystko możemy przejść do podpunktów.\\\\
(i)  Z jakim prawdopodobieństwem będziemy mieć więcej czarnych kul niż białych?\\
Wystarczy policzyć \(P(B>C)\). Zauważmy, że jako, że te zmienne mają ten sam rozkład to mamy symetrycznie \(P(B>C)=P(C>B)\). Stąd dostaniemy \(P(B>C)=\frac{1-P(B=C)}{2}\)
\[P(B=C)=\sum_{n=0}^\infty P(B=n \cap C=n) = \sum_{n=0}^\infty P(B=n) \cdot P(C=n)=\]
\[\sum_{n=0}^\infty \left(\frac{e^{-337}(337)^n}{n!}\right)^2\]
Stąd mamy odpowiedź:
\[P(B>C)=\frac{1-\sum_{n=0}^\infty \left(\frac{e^{-337}(337)^n}{n!}\right)^2}{2}\]
(ii) Jaki jest rozkład liczby czarnych kul przy założeniu, że nie uzyskaliśmy żadnej białej kuli?\\
Liczymy \(P(C|B=0)\) z definicji
\[P(C=k|B=0)=\frac{P(C=k \cap B=0)}{P(B=0)} = \frac{P(C=k) \cdot P( B=0)}{P(B=0)} = P(C=k)\]
Rozkład jest taki sam jak dla samej zmiennej C (warunek się skraca z niezależności), czyli:
\[P(C=k)=\frac{e^{-337}(337^k)}{k!}\]
(iii) Jaka jest oczekiwana liczba białych kul przy założeniu, że czarnych kul jest dwa razy więcej od białych?
Liczymy wartość oczekiwaną z definicji:
\[E(B|C=2B) = \sum_{k=0}^\infty kP(B=k|C=2B)=\sum_{k=0}^\infty \frac{kP(B=k\cap C=2k)}{P(C=2B)} =\]
\[ \frac{ \sum_{k=0}^\infty k P(B=k)\cdot P(C=2k)}{\sum_{j=0}^\infty P(B=j)\cdot P(C=2j)}\]
To już nam się tak ładnie nie skróci. Możemy jedynie podstawić wzór na rozkład Poissona:
\[E(B|C=2B) = \sum_{k}^\infty \frac{k \frac{e^{-337}(337^k)}{k!} \frac{e^{-337}(337^{2k})}{(2k)!}}{\sum_{j=0}^\infty \frac{e^{-337}(337^j)}{j!} \frac{e^{-337}(337^{2j})}{(2j)!}} = \]
\[ \frac{ \sum_{k=0}^\infty k \frac{(337^{3k})}{k!(2k)!}}{\sum_{j=0}^\infty \frac{(337^{3j})}{j!(2j)!}} \]
\subsection*{Odpowiedź}
\[\text{(i) } \frac{1-P(B=C)}{2}\]
\[\text{(ii) } P(C=k)\]
\[\text{(iii) }  \frac{ \sum_{k=0}^\infty k P(B=k)\cdot P(C=2k)}{\sum_{j=0}^\infty P(B=j)\cdot P(C=2j)}\]
